% Resumo expandido em portugues (Resolucao 114/2004 - CAMPG/UFRGS),
% recomendado quando o trabalho e escrito em ingles.
\chapter{Resumo Expandido}
\label{ap:resumo-expandido}

\begin{otherlanguage}{brazilian}

\begin{center}
\textbf{Avaliação do Impacto do Vazamento de Dados em Pipelines de Aprendizado de Máquina para Tarefas de Regressão}
\end{center}

\subsection*{Introdução}

Modelos de aprendizado de máquina (\emph{machine learning}, ML) são cada vez mais
usados para apoiar decisões em áreas como saúde, finanças e engenharia. Para que
os resultados reportados sejam confiáveis, a estimativa do desempenho de
generalização precisa ser obtida com dados que não participaram da construção do
modelo. O vazamento de dados (\emph{data leakage}) ocorre quando essa separação
é violada, isto é, quando informações que não estariam disponíveis no momento da
predição são usadas durante o desenvolvimento do modelo \cite{kaufman2012}. Uma
forma frequente e difícil de perceber é a aplicação de etapas de
pré-processamento ou de seleção de modelos sobre o conjunto de dados completo,
antes da divisão entre treino e teste \cite{kapoor2023}.

Estudos anteriores documentaram os efeitos do vazamento em domínios específicos,
como bioinformática e imagens médicas, e \citeonline{becker2023} avaliou
sistematicamente esses efeitos em tarefas de classificação. Para tarefas de
regressão, nas quais a variável alvo é contínua e as métricas de avaliação são
diferentes, não há evidência equivalente. O objetivo deste trabalho é quantificar
como o vazamento introduzido em quatro etapas de um pipeline (normalização,
imputação de valores faltantes, seleção de atributos e ajuste de
hiperparâmetros) distorce a avaliação de modelos de regressão, e investigar se
alguns algoritmos e conjuntos de dados são mais sensíveis a essas distorções.

\subsection*{Metodologia}

O estudo adota um desenho pareado. Para cada combinação de conjunto de dados,
algoritmo e cenário de vazamento, são executados dois pipelines que diferem
apenas na etapa estudada: no pipeline sem vazamento, todo componente que aprende
com os dados é ajustado somente nos dados de treino de cada partição; no pipeline
com vazamento, o componente estudado é ajustado no conjunto completo antes da
divisão. Ambos usam as mesmas partições e sementes aleatórias, e cada
configuração é repetida seis vezes.

Foram utilizados 21 conjuntos de dados de regressão do repositório PMLB
\cite{romano2022}, com 47 a 40.768 instâncias e 4 a 124 atributos, e cinco
algoritmos: regressão linear (LR), k vizinhos mais próximos (KNN), árvore de
decisão (DT), floresta aleatória (RF) e regressão por vetores de suporte (SVR).
Na normalização, foi usada a padronização $z$-score. Na imputação, valores
faltantes foram inseridos com um mecanismo MAR \cite{muzellec2020} em taxas de
5\% a 30\% e preenchidos pela média, pela mediana ou por KNN. Na seleção de
atributos, um filtro univariado baseado no teste $F$ manteve de 1\% a 20\% dos
atributos. No ajuste de hiperparâmetros, a busca em grade com validação cruzada
simples (que escolhe a configuração e reporta seu desempenho nos mesmos dados)
foi comparada à validação cruzada aninhada. No total, foram obtidas 11.340
comparações pareadas.

O efeito do vazamento foi medido pela diferença de $R^2$ entre os pipelines com e
sem vazamento e pela redução percentual do RMSE, ambas orientadas de modo que
valores positivos indicam desempenho inflado. As diferenças foram analisadas com
o teste de postos sinalizados de Wilcoxon sobre as médias por conjunto de dados,
com correção de Holm, com o teste de Friedman para comparar algoritmos e com a
correlação de Spearman para relacionar os efeitos às características dos
conjuntos de dados \cite{demsar2006}.

\subsection*{Resultados}

Os quatro cenários apresentaram comportamentos muito diferentes. O vazamento na
normalização e na imputação, que transmite apenas estatísticas dos atributos,
produziu diferenças desprezíveis (média de cerca de $2\times10^{-4}$ no $R^2$),
distribuídas de forma quase simétrica em torno de zero. A única exceção foi um
viés pequeno, mas estatisticamente significativo, para o SVR na imputação, mais
visível com o imputador KNN, que usa instâncias individuais para preencher os
valores faltantes.

O ajuste de hiperparâmetros sem validação cruzada aninhada foi o cenário mais
crítico: inflou significativamente as estimativas dos cinco algoritmos, com
diferenças de até 0,35 no $R^2$ e RMSE subestimado em até 54\%. O KNN foi o
algoritmo afetado de forma mais consistente, enquanto LR e SVR apresentaram os
casos individuais mais extremos. A seleção de atributos antes da divisão dos
dados não alterou 93,6\% dos pares, pois a seleção feita com e sem as instâncias
de teste costuma coincidir, mas produziu algumas das maiores distorções do
estudo: em um conjunto com 60 instâncias, o $R^2$ estimado para a regressão
linear passou de $-0{,}13$ para $0{,}49$.

Os maiores efeitos concentraram-se nos seis conjuntos de dados com menos de 1.000
instâncias. No ajuste de hiperparâmetros, a inflação mostrou forte correlação
negativa com o número de instâncias ($\rho = -0{,}82$) e positiva com a razão
entre atributos e instâncias ($\rho = 0{,}82$). Diferenças significativas entre
algoritmos foram encontradas apenas nesse cenário; os modelos baseados em
árvores, em especial a RF, foram os mais robustos de modo geral.

\subsection*{Conclusão}

Os resultados indicam que a gravidade do vazamento depende principalmente do tipo
de informação que vaza. Etapas que utilizam a variável alvo, direta ou
indiretamente, como a seleção supervisionada de atributos e o ajuste de
hiperparâmetros, são as fontes mais perigosas de vazamento em regressão, sobretudo
quando há poucos dados disponíveis. Recomenda-se, portanto, que todas as etapas
que aprendem com os dados sejam incluídas no pipeline validado e que o desempenho
de procedimentos de ajuste de hiperparâmetros seja estimado com validação cruzada
aninhada. Como trabalhos futuros, propõe-se investigar mecanismos de dados
faltantes mais severos, métodos de seleção de atributos do tipo \emph{wrapper} e
embutidos, a combinação de múltiplas fontes de vazamento e a variação controlada
do tamanho dos conjuntos de dados.

\end{otherlanguage}
